\title{Half-Cell Flux and an Exactly Solvable Period-Eight Phase in Signed Cycle Squares}
\author{Yicheng Zhao \and Jiachen Li}
\date{}
\maketitle

\begin{center}
\small
School of Mathematics and Statistics, Chongqing University of Posts and
Telecommunications, Chongqing, China\\[2pt]
Yicheng Zhao: \texttt{2023213805@cqupt.edu.cn},
ORCID \href{https://orcid.org/0009-0003-7618-1661}{0009-0003-7618-1661}\\
Jiachen Li: \texttt{2023213809@stu.cqupt.edu.cn},
ORCID \href{https://orcid.org/0009-0006-5119-3369}{0009-0006-5119-3369}
\end{center}

\begin{abstract}
We minimize the adjacency spectral radius over edge signings of the fixed
cycle square $C_n(1,2)$.  For an even-period Hamilton-gauge signing, we
characterize when alternating sign followed by half-cell translation is a
chiral involution: the local square word must be half-periodic with negative
half-cell flux.  This switching-invariant condition reduces a
$2m$-dimensional Bloch fiber to an $m$-dimensional squared problem.  Chirality
already occurs at period two, where the squared edge is $8$; period eight is
the smallest primitive period for which the edge is strictly smaller.  Its
unique sub-eight orbit has two antipodal defects.  For this phase we determine
all four squared dispersion branches and prove, for every $L\geq1$,
\[
 \rho(A_{8L,+}^{*})^2=4+\sqrt{10+2\sqrt5}.
\]
The negative-holonomy rings have an explicit radius from the same dispersion
law.  For every $L\geq4$, the positive-holonomy phase has smaller radius than
the natural twisted signing on $C_{8L}(1,2)$.
\end{abstract}

\noindent\textbf{Keywords:} signed graph; cycle square; spectral radius;
switching; Floquet decomposition; chiral symmetry

\medskip
\noindent\textbf{MSC 2020:} 05C50 (primary); 05C22, 15A18 (secondary)
